\subsection{Precision and Recall}
\label{sec:precision-recall}
\slidesrc{Terminology introduction: 43--51; Semantics (long version): 44--50}

Suppose a program searches a collection of documents and returns those it thinks are
relevant to a query. How good is it? Two questions matter, and they pull in opposite
directions:
\begin{itemize}[nosep]
  \item of what it \emph{returned}, how much is actually relevant? (\textbf{precision})
  \item of what is actually relevant, how much did it \emph{find}? (\textbf{recall})
\end{itemize}

\begin{center}
\begin{tikzpicture}
  \draw[thick, nbgray] (-3.8,-2.2) rectangle (4.8,2.3);
  \fill[nbblue!15]   (-0.7,0) circle (1.6);
  \fill[nborange!20] (1.7,0)  circle (1.6);
  \begin{scope}
    \clip (-0.7,0) circle (1.6);
    \fill[nbgreen!35] (1.7,0) circle (1.6);
  \end{scope}
  \draw[thick, nbblue]   (-0.7,0) circle (1.6);
  \draw[thick, nborange] (1.7,0)  circle (1.6);
  \node[nbblue, font=\small\bfseries]   at (-1.6,1.9) {relevant};
  \node[nborange, font=\small\bfseries] at (2.6,1.9)  {retrieved};
  \node[font=\footnotesize, align=center] at (-1.3,0) {relevant,\\missed\\(FN)};
  \node[font=\footnotesize\bfseries] at (0.5,0) {TP};
  \node[font=\footnotesize, align=center] at (2.3,0) {retrieved,\\not\\relevant\\(FP)};
  \node[font=\footnotesize, nbgray, anchor=south west] at (-3.8,-2.2) {corpus (all documents)};
  \node[font=\footnotesize, anchor=south east] at (4.8,-2.2) {neither (TN)};
\end{tikzpicture}
\captionof{figure}{The green overlap (TP) holds the documents that are both relevant
and retrieved. Precision looks at the orange circle (what fraction of it is green),
recall looks at the blue circle (what fraction of it is green).}
\end{center}

\begin{nbdef}{Precision and recall}
\begin{align*}
  \text{precision} &= \frac{\abs{\{\text{relevant documents}\} \cap \{\text{retrieved documents}\}}}
                           {\abs{\{\text{retrieved documents}\}}} \\[4pt]
  \text{recall}    &= \frac{\abs{\{\text{relevant documents}\} \cap \{\text{retrieved documents}\}}}
                           {\abs{\{\text{relevant documents}\}}}
\end{align*}
\textbf{Precision} (also called \emph{positive predictive value}) is the fraction of
retrieved documents that are relevant to the query. \textbf{Recall} (also called
\emph{sensitivity}) is the fraction of relevant documents that are successfully
retrieved. Both are between 0 and 1; $\abs{S}$ denotes the number of elements of the
set $S$.
\end{nbdef}

The numerator is the \emph{same} in both formulas: the documents that are both relevant
and retrieved. Only the denominator changes: precision divides by what the system
\emph{returned}, recall by what it \emph{should have returned}.

\subsubsection{The same thing with a confusion matrix}
\notinslides

Each document falls in one of four cases:

\begin{center}
\begin{tabular}{@{}l c c@{}}
\toprule
 & \textbf{relevant} & \textbf{not relevant} \\
\midrule
\textbf{retrieved}     & true positive (TP)  & false positive (FP) \\
\textbf{not retrieved} & false negative (FN) & true negative (TN) \\
\bottomrule
\end{tabular}
\end{center}

\[
  \text{precision} = \frac{TP}{TP + FP}, \qquad
  \text{recall}    = \frac{TP}{TP + FN}.
\]
Neither formula uses TN: the (usually huge) number of irrelevant documents that were
correctly left out does not count.

\subsubsection{Four worked examples}
\slidesrc{Terminology introduction: 44--51}

The setting is the same for all four examples. You implement an application that
retrieves the documents containing one specific word (``cat'') and its inflectional
forms (``cats'', see Section~\ref{sec:morphology}) in a corpus of \textbf{1000}
documents. Exactly \textbf{24} documents contain ``cat'': these are the \textbf{relevant}
documents.

\begin{nbexample}{1: a typical system}
The application retrieves \textbf{36} documents; \textbf{16} of them contain ``cat'' and
20 do not. So $TP=16$, $FP=20$, $FN = 24-16 = 8$, $TN = 1000-16-20-8 = 956$.
\begin{align*}
  \text{precision} &= \frac{16}{36} = \frac{4}{9} \approx 0.444
     \reason{$TP/(TP+FP) = 16/(16+20)$} \\
  \text{recall}    &= \frac{16}{24} = \frac{2}{3} \approx 0.667
     \reason{$TP/(TP+FN) = 16/(16+8)$}
\end{align*}
Fewer than half of the results are right, and a third of the cat documents are missed.
\end{nbexample}

\begin{nbexample}{2: return everything}
The application retrieves \textbf{all 1000} documents; \textbf{24} of them contain
``cat'' and 976 do not. So $TP = 24$, $FP = 976$, $FN = 0$, $TN = 0$.
\begin{align*}
  \text{precision} &= \frac{24}{1000} = 0.024 \\
  \text{recall}    &= \frac{24}{24} = 1
\end{align*}
Perfect recall, for free: a system that returns everything never misses anything. Its
precision is terrible.
\end{nbexample}

\begin{nbwarning}[a typo in the slides]
Example~2 on the slides says the retrieved documents ``out of which 24 contain the word
cat, and 1986 do not contain it''. Since only 1000 documents are retrieved, the correct
number is $1000 - 24 = 976$. The precision and recall values on the slides are correct.
\end{nbwarning}

\begin{nbexample}{3: play it safe}
The application retrieves only \textbf{2} documents, and both contain ``cat''. So
$TP=2$, $FP=0$, $FN=22$, $TN=976$.
\begin{align*}
  \text{precision} &= \frac{2}{2} = 1 \\
  \text{recall}    &= \frac{2}{24} = \frac{1}{12} \approx 0.083
\end{align*}
Perfect precision, almost for free: return only the one or two results you are
completely sure about. Its recall is terrible.
\end{nbexample}

\begin{nbexample}{4: the perfect system}
The application retrieves \textbf{24} documents, all containing ``cat''. So $TP=24$,
$FP=0$, $FN=0$, $TN=976$.
\[
  \text{precision} = \frac{24}{24} = 1, \qquad \text{recall} = \frac{24}{24} = 1.
\]
\end{nbexample}

\begin{center}
\begin{tikzpicture}
\begin{axis}[ybar, bar width=13pt, width=\linewidth, height=6.5cm,
    ymin=0, ymax=1.15, ytick={0,0.25,0.5,0.75,1},
    symbolic x coords={Ex.\ 1, Ex.\ 2, Ex.\ 3, Ex.\ 4}, xtick=data,
    tick label style={font=\scriptsize}, label style={font=\small},
    legend style={font=\small, at={(0.5,1.03)}, anchor=south, legend columns=3, draw=none},
    enlarge x limits=0.18,
    nodes near coords, every node near coord/.append style={font=\scriptsize, /pgf/number format/fixed, /pgf/number format/precision=2}]
  \addplot[fill=nborange!70, draw=nborange] coordinates {(Ex.\ 1,0.444) (Ex.\ 2,0.024) (Ex.\ 3,1) (Ex.\ 4,1)};
  \addplot[fill=nbblue!60, draw=nbblue]     coordinates {(Ex.\ 1,0.667) (Ex.\ 2,1) (Ex.\ 3,0.083) (Ex.\ 4,1)};
  \addplot[fill=nbgreen!60, draw=nbgreen]   coordinates {(Ex.\ 1,0.533) (Ex.\ 2,0.047) (Ex.\ 3,0.154) (Ex.\ 4,1)};
  \legend{precision, recall, $F_1$}
\end{axis}
\end{tikzpicture}
\captionof{figure}{Precision, recall and $F_1$ for the four examples. Examples 2 and 3
each score perfectly on one measure with a trivial strategy; only $F_1$ exposes them.}
\end{center}

\begin{nbintuition}[the trade-off]
Precision and recall must always be read \emph{together}. Returning more documents
tends to raise recall and lower precision (Example~2 is the extreme); returning fewer
tends to raise precision and lower recall (Example~3). A single number that is high only
when \emph{both} are high is the $F_1$ score below.
\end{nbintuition}

\begin{nbdef}{$F_1$ score (not on the slides)}
The \textbf{$F_1$ score} is the harmonic mean of precision $P$ and recall $R$:
\[
  F_1 = \frac{2PR}{P + R} = \frac{2\,TP}{2\,TP + FP + FN}.
\]
The harmonic mean is dominated by the smaller of the two values, so a system cannot get a
good $F_1$ by maximising only one of them.
\end{nbdef}

For the four examples:
\begin{align*}
  \text{Ex.\ 1:}\quad F_1 &= \frac{2\cdot 16}{2\cdot 16 + 20 + 8} = \frac{32}{60} \approx 0.533 \\
  \text{Ex.\ 2:}\quad F_1 &= \frac{2\cdot 24}{2\cdot 24 + 976 + 0} = \frac{48}{1024} \approx 0.047 \\
  \text{Ex.\ 3:}\quad F_1 &= \frac{2\cdot 2}{2\cdot 2 + 0 + 22} = \frac{4}{26} \approx 0.154 \\
  \text{Ex.\ 4:}\quad F_1 &= \frac{2\cdot 24}{2\cdot 24 + 0 + 0} = 1
\end{align*}

\subsubsection{Another example: five algorithms}
\slidesrc{Semantics (long version): 44--50}

The long version of the semantics slides repeats the definitions with a second example.
A query asks for the documents about topic $T$ in a corpus $C$. The correct answer is
known: $D_1, D_2, D_4, D_5, D_6, D_7, D_8$ deal with $T$ (7 relevant documents) and no
other document does.

\begin{center}
\footnotesize
\begin{tabular}{@{}l l c c c@{}}
\toprule
\textbf{Algorithm} & \textbf{Returns} & $\abs{\text{rel}\cap\text{ret}}$ & \textbf{Precision} & \textbf{Recall} \\
\midrule
A1 & $D_1, D_2, D_3, D_4$            & 3 & $3/4 = 75\%$       & $3/7 \approx 43\%$ \\
A2 & $D_1, \dots, D_8$               & 7 & $7/8 = 87.5\%$     & $7/7 = 100\%$ \\
A3 & $D_1$ only                      & 1 & $1/1 = 100\%$      & $1/7 \approx 14\%$ \\
A4 & all and only the relevant & 7 & $7/7 = 100\%$      & $7/7 = 100\%$ \\
A5 & 1000 docs, incl.\ the 7 relevant & 7 & $7/1000 = 0.7\%$ & $7/7 = 100\%$ \\
\bottomrule
\end{tabular}
\end{center}

The same lessons as before: A3 buys perfect precision by returning almost nothing, A5
buys perfect recall by returning almost everything; only A4 is perfect on both. (On the
slide for A2 the formulas are labelled ``precision(A1)'' and ``recall(A1)'': a
copy-paste slip, the values are those of A2.)

\subsubsection{Exercises}
\notinslides

\begin{nbexercise}{1}{morphemes}
Split each word into morphemes and classify each morpheme as free or bound, and each
bound morpheme as derivational or inflectional: \emph{unhappiness}, \emph{walked},
\emph{teachers}, \emph{rewritten}. Which of \emph{bookshelf}, \emph{bookish},
\emph{books} is a compound?
\end{nbexercise}
\begin{nbsolution}
\begin{itemize}[nosep]
  \item \emph{un-happi-ness}: \emph{un-} bound, derivational (negates the meaning);
        \emph{happy} free root; \emph{-ness} bound, derivational (adjective $\to$ noun).
  \item \emph{walk-ed}: \emph{walk} free root; \emph{-ed} bound, inflectional (past tense).
  \item \emph{teach-er-s}: \emph{teach} free root; \emph{-er} bound, derivational
        (verb $\to$ noun, ``one who teaches''); \emph{-s} bound, inflectional (plural).
  \item \emph{re-writ-ten}: \emph{re-} bound, derivational (``again''); \emph{write}
        free root; \emph{-en} bound, inflectional (past participle).
\end{itemize}
Only \emph{bookshelf} is a compound (\emph{book} + \emph{shelf}, both free).
\emph{Bookish} is derived (\emph{-ish} is bound), \emph{books} is inflected.

\emph{Lesson}: derivational morphemes sit closer to the root than inflectional ones
(\emph{teach-er-s}, not \emph{teach-s-er}).
\end{nbsolution}

\begin{nbexercise}{2}{precision and recall}
A corpus has 500 documents, 40 of which are relevant to a query. A search engine
returns 50 documents, 30 of which are relevant. Compute TP, FP, FN, TN, precision,
recall and $F_1$. What happens to precision and recall if the engine also returns 50 more
documents, none of them relevant?
\end{nbexercise}
\begin{nbsolution}
$TP = 30$, $FP = 50 - 30 = 20$, $FN = 40 - 30 = 10$, $TN = 500 - 30 - 20 - 10 = 440$.
\begin{align*}
  \text{precision} &= \frac{30}{50} = 0.6, \qquad
  \text{recall} = \frac{30}{40} = 0.75, \\
  F_1 &= \frac{2\cdot 30}{2\cdot 30 + 20 + 10} = \frac{60}{90} \approx 0.667.
\end{align*}
With 50 more irrelevant documents: $TP = 30$, $FP = 70$, $FN = 10$, so
precision $= 30/100 = 0.3$ (halved) while recall stays $30/40 = 0.75$, and
$F_1 = 60/(60+70+10) = 60/140 \approx 0.429$.

\emph{Lesson}: adding wrong results never changes recall but always lowers precision.
\end{nbsolution}

\begin{nbsummary}[terminology]
\begin{itemize}[nosep]
  \item \textbf{Morpheme}: smallest meaningful unit; \textbf{free} (stands alone:
        \emph{dog}) or \textbf{bound} (\emph{un-}, \emph{-ness}, \emph{-s}).
  \item Bound morphemes: \textbf{derivational} (new word: meaning or part of speech
        changes, \emph{happy} $\to$ \emph{happiness}) or \textbf{inflectional} (new form
        of the same word: tense, number, \ldots; \emph{dog} $\to$ \emph{dogs}).
  \item \textbf{Word}: smallest element uttered in isolation with meaning; one or more
        morphemes. \textbf{Compound}: word made of $\geq 2$ \emph{free} morphemes.
  \item \textbf{Corpus}: large structured set of texts. \textbf{Genre of a corpus}: from
        text-external features (medium, function, format). \textbf{Collocation}: words
        co-occurring more often than by chance (\emph{fast food}, not \emph{quick food}).
  \item $\text{precision} = TP/(TP+FP)$: of what I returned, how much is right;
        $\text{recall} = TP/(TP+FN)$: of what is right, how much I returned;
        $F_1 = 2PR/(P+R)$. Returning everything gives recall 1; returning very few
        sure hits gives precision 1.
\end{itemize}
\end{nbsummary}
